\chapter{Technology Review}
This chapter discusses the different technologies used in throughout the project. It discusses the the advantages and disadvantages of each technology and why certain technologies were used over others.

\section{Overview}
This project is comprised of React as the frontend, Flask as the server, MongoDB as the database, and ... to host the application. Throughout the project, the following technologies were also used and tested before the above was ultimately chosen:
\begin{itemize}
    \item Angular/Ionic
    \item Firebase
    \item Django
    \item MySQL
    \item Amazon Web Services
\end{itemize}

\section{Main Technologies}
This section will discuss the main technologies currently in use in the web application. The preceding section will discuss other technologies tried but ultimately weren't incorporated. 

\newpage

\subsection{React}
\par
\medskip
\begin{center}
    \includegraphics[width=8cm,height=3.3cm,keepaspectratio]{images/react}
\end{center}
React (also known as React.js or ReactJS) is a JavaScript library for building user interfaces. It is maintained by Facebook and a community of individual developers and companies.

React can be used as a base in the development of single-page or mobile applications. However React is only concerned with rendering data to the DOM and so creating React applications usually requires the use of additional libraries for state management and routing. Redux and React Router are respective examples of such libraries. 

\subsubsection{Advantages}
React has many advantages, several of which apply to this project:

% https://www.altexsoft.com/blog/engineering/the-good-and-the-bad-of-reactjs-and-react-native/
\paragraph{Virtual DOM in ReactJS makes user experience better and developer’s work faster}

\paragraph{Redux: convenient state container}

\paragraph{Wide React and Redux toolset}

\subsubsection{Disadvantages}


\subsection{Flask}
\par
\medskip
\begin{center}
    \includegraphics[width=8cm,height=3.3cm,keepaspectratio]{images/flask}
\end{center}

\subsubsection{Advantages}

\subsubsection{Disadvantages}


\subsection{MongoDB}
\par
\medskip
\begin{center}
    \includegraphics[width=8cm,height=3.3cm,keepaspectratio]{images/mongodb}
\end{center}

\subsubsection{Advantages}

\subsubsection{Disadvantages}

\subsection{PythonAnywhere}
\par
\medskip
\begin{center}
    \includegraphics[width=8cm,height=3.3cm,keepaspectratio]{images/pythonanywhere}
\end{center}
PythonAnywhere is an online integrated development environment (IDE) and web hosting service (Platform as a service) based on the Python programming language. it provides in-browser access to server-based Python and Bash command-line interfaces, along with a code editor with syntax highlighting. Program files can be transferred to and from the service using the user's browser. Web applications hosted by the service can be written using any WSGI-based application framework. For this project, the user authentication side of things was initially handled entirely by Firebase. However, as suggested by my supervisor, I decided to implement registering a user, logging a user into the application etc. myself. The functionality for this was written in Python and could be accessed by the application by using a proxy to delegate any requests made by the application to the Flask server. However, to enable the functionality to be accessed from anywhere (and, in turn, remove the need to run the Flask server alongside the application everytime), I decided to use PythonAnywhere. EXPLAIN

\subsubsection{Advantages}
There a number of advantages to using PythonAnywhere:

\paragraph{Always running}
Even on a free tier account, PythonAnywhere never sleeps (compared to something like Heroku). This means real time services are viable with PythonAnywhere.

\paragraph{Fully configured}
PythonAnywhere gives you a fully configured Python environment...

\paragraph{Free}
PythonAnywhere is free to use. On a free tier account, web applications stay running 24/7 for 3 months. After 3 months, it will shut down. However, one needs only to start it up again.

\subsubsection{Disadvantages}

\paragraph{Python-only on the server side}
You are free to use JavaScript in your web pages and so on, but you can't use Rails or Node on the server side of things.

\paragraph{No WebSocket support}
